\begin{longtable}{@{}p{1.5cm}p{3.0cm}p{8.55cm}@{}}
\caption{Kebutuhan nonfungsional sistem.}
\label{tab:kebutuhan-nonfungsional} \\
\toprule
\textbf{Kode} & \textbf{Atribut} & \textbf{Kebutuhan Nonfungsional} \\
\midrule
\endfirsthead

\multicolumn{3}{@{}l}{\small\emph{Tabel \thetable{} (lanjutan)}} \\
\toprule
\textbf{Kode} & \textbf{Atribut} & \textbf{Kebutuhan Nonfungsional} \\
\midrule
\endhead

\bottomrule
\endlastfoot

KNF-01 & Keterterapan & Sistem berjalan pada laptop kelas konsumen tanpa akselerator grafis. \\[2pt]

KNF-02 & Keterjelasan & Prediksi dapat dijelaskan melalui kontribusi fitur, dan rekomendasi dapat ditelusuri ke dokumen sumbernya. \\[2pt]

KNF-03 & Keamanan konten & Rekomendasi dibumikan pada dokumen, dengan parameter pembangkitan yang konservatif untuk menekan pemunculan angka dosis atau ambang yang tidak berdasar. \\[2pt]

KNF-04 & Ketahanan & Tersedia mekanisme cadangan yang terkendali pada setiap lapisan, mencakup lapisan data, penelusuran, dan pembangkitan. \\[2pt]

KNF-05 & Bahasa & Keluaran sistem menggunakan bahasa Indonesia. \\[2pt]

KNF-06 & Keamanan klinis & Sistem berperan sebagai pendukung keputusan dengan dokter sebagai penentu akhir, dan bukan sebagai pengganti penilaian klinis. \\[2pt]

KNF-07 & Keterpeliharaan & Terdapat pemisahan antara logika domain dan antarmuka pengguna. \\[2pt]

KNF-08 & Keterlacakan sitasi & Setiap potongan dokumen menyimpan metadata identitas dokumen, nomor halaman berkas, dan nomor halaman cetak, sehingga rekomendasi dapat diverifikasi sampai ke halaman sumbernya. \\[2pt]

KNF-09 & Biaya dan keterpindahan & Sistem memanfaatkan komponen bersumber terbuka atau berjenjang gratis, tanpa memerlukan server khusus. \\[2pt]

KNF-10 & Waktu tanggap & Sistem menyajikan prediksi dan rekomendasi dalam rentang waktu yang memungkinkan penggunaannya pada saat konsultasi berlangsung. Nilai acuan ditetapkan berdasarkan pengukuran pada Bab~\ref{chap:evaluasi}. \\

\end{longtable}
